%% frag_results.tex
%% Replaces \subsection{Retrieval Quality} in Section 4 and adds the
%% localization, ablation, and scaled-evaluation subsections after it.

\subsection{Two Metrics, Two Orderings}\label{sec:two-metrics}

We first report retrieval quality under the score used in the conference
version of this work, and then under standard information-retrieval metrics.
The two disagree, and the disagreement is the entry point to this section's
main result.

The bespoke score is a keyword- and source-overlap measure between the question
and the retrieved set, computed by the Evaluator agent
(Section~\ref{sec:method}). Under it, the multi-agent architecture improved on
the single-agent baseline in all five question categories, from 0.680 to 0.798
overall --- a 17.2\% relative improvement, with a paired $t$-test over the 50
evaluation queries yielding $t(49)=2.32$, $p=0.020$.

We then re-scored the same pipeline with nDCG@$k$, Recall@$k$, and MRR. Graded
relevance labels were pooled across systems: every document retrieved by any
system for a question was judged once, and the resulting label was reused for
every system that retrieved it, so no system is scored against a private set of
judgments. Table~\ref{tab:inversion} reports the result on the ten
ground-truth validation questions.

\begin{table}[t]
\caption{The same two systems under the paper's own retrieval-quality score and
under standard IR metrics. Bespoke score: $n=50$. IR metrics: $n=10$, pooled
graded relevance, judge \texttt{llama3.1}.}
\label{tab:inversion}
\begin{tabular}{lrrrr}
\toprule
System & bespoke score & nDCG@1 & nDCG@3 & MRR \\
\midrule
Single-agent baseline & 0.680 & \textbf{0.800} & \textbf{0.765} & \textbf{0.800} \\
Multi-agent           & \textbf{0.798} & 0.133 & 0.145 & 0.250 \\
\bottomrule
\end{tabular}
\end{table}

The ordering reverses. A 17.2\% advantage on one metric coexists with roughly a
five-fold deficit on another, over the same questions, the same live corpus, and
the same retrieved documents. Because the bespoke score rewards overlap between
the question and the retrieved set without reference to whether a retrieved
document actually answers the question, it can be maximized by retrieving many
topically adjacent documents. Standard IR metrics, scored against judged
relevance, cannot be.

\subsection{Localizing the Loss: Ranking or Retrieval?}\label{sec:localizing}

A deficit of this size has two candidate explanations. Either the pipeline
retrieves appropriate documents and orders them badly, or it does not retrieve
them. These call for different fixes, and they are indistinguishable from the
final top-$k$ alone.

The vocabulary-mismatch framing that motivates query rewriting in the literature
invites the first reading, and it is the one we pursued first. The retriever's
ordering was a substring boost computed over the first four tokens of the
rewritten query, which after a rewrite are typically filler
(Section~\ref{sec:rerank}). We replaced it with three progressively stronger
rerankers scored against the user's original question.
Table~\ref{tab:ranking-ablation} reports the outcome.

\begin{table}[t]
\caption{Ranking ablation on the multi-agent arm ($n=10$). Successively stronger
rerankers raise MRR monotonically while nDCG@3 remains flat --- the signature of
a candidate set that does not contain the answer.}
\label{tab:ranking-ablation}
\begin{tabular}{llrrr}
\toprule
Configuration & Ranking signal & nDCG@3 & MRR & Recall@5 \\
\midrule
Published & first-4-token substring boost & 0.145 & 0.250 & 0.145 \\
BM25       & Okapi BM25 over the pool      & 0.212 & 0.500 & 0.203 \\
Embedding  & cosine over local embeddings  & 0.188 & 0.625 & 0.327 \\
\midrule
Embedding + union fetch & \emph{as above, wider pool} & \textbf{0.973} & \textbf{1.000} & \textbf{0.933} \\
\bottomrule
\end{tabular}
\end{table}

MRR rises monotonically from 0.250 to 0.625, confirming that ranking was
genuinely defective: better ranking does place better documents first among
those available. But nDCG@3 stays near 0.19 throughout. A ranking intervention
that improves ranking without closing the gap is evidence that the gap is not a
ranking problem.

We therefore stopped inferring from scores and counted documents directly. For
each question we took the documents the single-agent baseline retrieved that the
judge had labelled relevant, and asked whether each appeared anywhere in the
multi-agent pipeline's retrieved set:

\begin{quote}
Of the 23 relevant documents the single-agent baseline retrieved across the
evaluation set, \textbf{22 were never retrieved by the multi-agent pipeline at
all.}
\end{quote}

The cause is a wiring property rather than a modelling one. Every retrieval
endpoint was being issued the rewritten query \emph{instead of} the user's
question. The rewrite is a lossy paraphrase --- it is designed to be, since its
purpose is to substitute domain vocabulary for user vocabulary --- and the
documents matching the user's own phrasing were consequently never candidates.
No reranker can promote a document that was never retrieved, which is precisely
what Table~\ref{tab:ranking-ablation} shows.

The remedy follows directly and is stated in Section~\ref{sec:union-fetch}: issue
both phrasings and union the candidate pools. The final row of
Table~\ref{tab:ranking-ablation} gives its effect --- nDCG@3 from 0.188 to 0.973
in a single change, with no modification to ranking, generation, or the
evaluator. Notably, the multi-agent pipeline had all along been retrieving
relevant documents the baseline missed; its candidate set was not worse so much
as \emph{different and smaller}, which is why unioning rather than replacing
recovers the loss.

\subsection{Evaluation at Scale}\label{sec:scaled-eval}

The preceding measurements use ten ground-truth questions and are diagnostic
rather than conclusive. To test whether the repaired pipeline holds up, we
evaluated on a 100-question stratified subset of the 300-question benchmark
(Section~\ref{sec:method}), balanced across all five question categories at 20
each and covering 21 of the 24 software ecosystems. Ninety-six questions
completed across all three systems and are reported here.

Three systems were run with the synthesis model held constant at Llama 3.1 8B,
which isolates architecture from model:

\begin{itemize}
  \item \textbf{marag} --- the multi-agent pipeline as described, whose answer is
  a template assembled by the Evaluator agent from the retrieved documents.
  \item \textbf{marag\_llm} --- identical retrieval, but the answer is
  synthesized in prose through the same prompt the baseline uses. This arm
  exists solely so that answer-quality scores are comparable.
  \item \textbf{single\_agent} --- the baseline: raw query, one retrieval pass,
  one synthesis call.
\end{itemize}

\begin{table}[t]
\caption{Retrieval quality on the 100-question stratified benchmark subset
($n=96$), pooled graded relevance, judge \texttt{llama3.1}. Brackets give
bootstrap 95\% confidence intervals. The two retrieval arms are identical by
construction; they differ only in answer synthesis.}
\label{tab:scaled-retrieval}
\begin{tabular}{lrrrr}
\toprule
System & nDCG@3 & nDCG@5 & MRR & Recall@5 \\
\midrule
marag         & 0.865 [0.80, 0.92] & 0.867 [0.81, 0.92] & 0.892 [0.83, 0.95] & 0.866 [0.80, 0.92] \\
marag\_llm    & 0.865 [0.80, 0.92] & 0.867 [0.81, 0.92] & 0.892 [0.83, 0.95] & 0.866 [0.80, 0.92] \\
single\_agent & 0.869 [0.81, 0.92] & 0.876 [0.82, 0.93] & 0.895 [0.83, 0.95] & 0.875 [0.82, 0.93] \\
\bottomrule
\end{tabular}
\end{table}

Table~\ref{tab:scaled-retrieval} reports retrieval quality. The repaired
multi-agent pipeline and the single-agent baseline are indistinguishable: every
difference is well within overlapping confidence intervals, and the baseline is
nominally ahead on all four metrics. The fetch-time diagnostic of
Section~\ref{sec:localizing}, recomputed on this larger set, falls from 96\% (22
of 23) to \textbf{14.7\%} (23 of 156): the union remedy holds at scale, and the
residual is the ordinary variation of two systems retrieving from live sources.

\begin{table}[t]
\caption{Answer quality and cost ($n=96$), LLM-judged. \texttt{marag} and
\texttt{marag\_llm} share a retrieval pipeline and differ only in how the answer
is produced. Latency is the per-question median in seconds.}
\label{tab:scaled-answer}
\begin{tabular}{lrrr}
\toprule
System & Faithfulness & Answer relevance & Latency (s) \\
\midrule
marag (template answer) & 0.702 & 0.977 & 19.6 \\
marag\_llm (prose answer) & 0.924 & 0.952 & 21.1 \\
single\_agent & 0.929 & 0.973 & 9.5 \\
\bottomrule
\end{tabular}
\end{table}

Table~\ref{tab:scaled-answer} carries a methodological result worth stating
separately. Comparing \texttt{marag} against \texttt{single\_agent} on
faithfulness suggests a substantial deficit (0.702 versus 0.929). Comparing
\texttt{marag\_llm} against \texttt{single\_agent} --- the same retrieval, the
same documents, differing only in whether the answer is a structured template or
model prose --- gives 0.924 versus 0.929. The apparent deficit was almost
entirely an artifact of answer \emph{format}: an LLM judge scoring a structured
template against fluent prose is partly scoring fluency. Any evaluation
comparing an agentic system that emits structured output against a baseline that
emits prose is exposed to this, and we would not have detected it without an arm
constructed specifically to remove it.

On cost, the multi-agent pipeline takes approximately twice the median
per-question latency of the baseline (19.6\,s versus 9.5\,s). This reflects the
rewriter's generation call, the evaluator pass, and the additional retrieval
calls the union incurs.

\subsection{Summary of Findings}\label{sec:results-summary}

\begin{enumerate}
  \item The bespoke retrieval-quality score and standard IR metrics order the
  two architectures oppositely on the same data (Table~\ref{tab:inversion}).
  \item The multi-agent pipeline's deficit under IR metrics was caused by
  retrieval, not ranking: 22 of 23 relevant documents were never retrieved,
  and three reranker generations could not recover them
  (Table~\ref{tab:ranking-ablation}).
  \item Issuing both the original and rewritten phrasings and unioning the
  candidate pools removes the deficit, and the remedy holds on a
  100-question multi-ecosystem benchmark, where the fetch-time loss falls from
  96\% to 14.7\%.
  \item With both defects repaired and the synthesis model held constant, the
  multi-agent architecture \emph{matches} the single-agent baseline on retrieval
  and on faithfulness, at roughly twice the latency
  (Tables~\ref{tab:scaled-retrieval} and~\ref{tab:scaled-answer}). We do not
  find evidence that decomposition improves retrieval quality in this domain.
  \item A substantial apparent faithfulness gap between the two architectures
  was an artifact of answer format rather than answer content.
\end{enumerate}
